\section{A dimension-explicit Ricci reaction bound}\label{sec:reaction}

The constant in Proposition~\ref{prop:full-curvature} may grow arbitrarily
with the dimension. We next obtain a different estimate whose leading
dimension dependence is controlled.
Only the latter estimate will be used to choose the dimension. The
cap estimate will enter afterward, with its flow-dependent constant,
through an independently chosen small parameter.

\subsection{A one-dimensional interior estimate}

For \(R>0\), let \(P_R=(-R,R)\times[-R^2,0]\), with space coordinate
\(y\) and time coordinate \(\rho\).

\begin{lemma}\label{lem:heat-perturbation}
There is \(\epsilon_0>0\) with the following property. If a smooth
function \(H\) on \(P_2\) satisfies
\[
H_\rho=aH_{yy}+bH_y+F,\qquad
|a-1|\le\epsilon_0,\qquad |H|+|b|+|F|\le K,
\]
then \(|H_y|\le C(K)\) on \(P_1\), including the terminal time by
continuity from below. No bounds on derivatives of the coefficients are
required. The same conclusion holds after any fixed rescaling of the
cylinders.
\end{lemma}

\begin{proof}
Fix an exponent \(s>3\). The constant-coefficient heat estimate is
\begin{equation}\label{eq:heat-Lp}
\|V_{yy}\|_{L^s}+\|V_\rho\|_{L^s}
\le C_s\|V_\rho-V_{yy}\|_{L^s}
\end{equation}
for smooth Sobolev functions compactly supported in space and vanishing in the distant
past, with norms over times up to zero. This is the classical parabolic
Calder\'on--Zygmund estimate for the heat operator; its constants here are
one-dimensional; see~\cite[Theorem~5.5]{Krylov2006}, with time reversed.
It follows equivalently from the \(L^s\) boundedness of
the second spatial derivative of the causal heat potential. The
half-infinite time version follows by extending its forcing to future
times by zero. The function itself need not vanish at time zero.

For \(1\le R_1<R_2<2\), put \(d_R=R_2-R_1\) and choose a cutoff
\(\zeta\) equal to one on \(P_{R_1}\), supported spatially and in the
past inside \(P_{R_2}\), with
\[
|\zeta_y|\le C d_R^{-1},\qquad
|\zeta_{yy}|+|\zeta_\rho|\le C d_R^{-2}.
\]
Writing \(V=\zeta H\), direct expansion gives
\begin{equation}\label{eq:cutoff-heat}
\begin{split}
(\partial_\rho-\partial_{yy})V={}&(a-1)V_{yy}+bV_y+\zeta F\\
&+(\zeta_\rho-a\zeta_{yy}-b\zeta_y)H-2a\zeta_yH_y.
\end{split}
\end{equation}
Choose \(\epsilon_0\) so that \(C_s\epsilon_0<1/2\) and absorb the first
term using~\eqref{eq:heat-Lp}. If
\(E(R)=\|H_{yy}\|_{L^s(P_R)}+\|H_\rho\|_{L^s(P_R)}\), the remaining
terms yield
\[
E(R_1)\le C(K)d_R^{-1}\|H_y\|_{L^s(P_{R_2})}+C(K)d_R^{-2}.
\]
Spatial interpolation on intervals of radius between one and two gives
\[
\|H_y\|_{L^s(P_R)}\le h\|H_{yy}\|_{L^s(P_R)}
       +C h^{-1}\|H\|_{L^s(P_R)}
\]
for sufficiently small \(h>0\), without boundary conditions on \(H\).
Taking \(h\) to be a sufficiently small multiple of \(d_R\) gives
\begin{equation}\label{eq:interior-iteration}
E(R_1)\le\tfrac1{16}E(R_2)+C(K)d_R^{-2}.
\end{equation}
Iterate with \(R_j=7/4-2^{-j}/4\). The error terms are summable since
their factors grow as \(4^j\) whereas the iteration contributes
\(16^{-j}\). The remainder tends to zero because
\(E(R_j)\le E(7/4)<\infty\) for each individual smooth function. This
finiteness is not a presumed uniform bound. Hence \(E(3/2)\le C(K)\).
Interpolation also controls the first derivative norm. Parabolic Sobolev
embedding~\cite[Chapter~II, Section~3, Lemma~3.3]{LSU1968},
with \(s>1+2\), bounds the spatial first derivative on
\(P_1\). This proves the claim.
\end{proof}

\subsection{Cylinders moving with the radial drift}

\begin{lemma}\label{lem:largeq-derivatives}
For all sufficiently large \(q\), and sufficiently late times,
\begin{equation}\label{eq:largeq-derivatives}
|rz_r|\le Cq^{-1/2},\qquad r^2|u_{rr}|\le Cq^{1/2}
\quad\left(0<r\le\tfrac18\sqrt\delta\right),
\end{equation}
where \(C\) is independent of both \(q\) and \(k\).
\end{lemma}

\begin{proof}
Fix \((r_c,t_c)\) in the indicated region and set
\begin{equation}\label{eq:moving-cylinder}
h_c=\frac{r_c}{\sqrt q},\qquad
t=t_c+h_c^2\rho,\qquad
r=R_c(t)+h_cy,\qquad
R_c(t)=\sqrt{r_c^2+2(q-1)(t_c-t)}.
\end{equation}
On \(P_2\), \(1\le R_c/r_c\le3\). For \(q\ge16\),
\begin{equation}\label{eq:cylinder-inclusion}
\frac12r_c\le r\le\frac72r_c,
\qquad r^2\le\frac{49}{256}\delta_c<\delta(t),
\qquad \delta_c=T-t_c.
\end{equation}
Moreover its earliest time is at least
\(t_c-\delta_c/(16q)\). Thus a single sufficiently late threshold for
\(t_c\), for the chosen flow, puts every such cylinder inside the
region of Lemma~\ref{lem:uniform}.

The moving center in~\eqref{eq:moving-cylinder} follows the radial drift.
Figure~\ref{fig:drift} depicts the resulting cylinder. A cylinder centered
at a fixed radius would leave an uncontrolled coefficient of size
\(\sqrt q\); the movement cancels this coefficient before the
one-dimensional estimate is applied.

\begin{figure}[htbp]
\centering
\begin{tikzpicture}[x=1cm,y=1cm,>=Latex,font=\small]
  \draw[->] (0,0)--(10.2,0) node[right] {$r$};
  \draw[->] (0,0)--(0,4.5) node[above] {$t$};
  \fill[blue!10] (7.2,.5) .. controls (6.3,1.6) and (5.6,2.7) .. (4.2,3.8)
      --(6.2,3.8) .. controls (7.6,2.7) and (8.3,1.6) .. (9.2,.5)--cycle;
  \draw[blue!65!black,thick] (7.2,.5) .. controls (6.3,1.6) and (5.6,2.7) .. (4.2,3.8);
  \draw[blue!65!black,thick] (9.2,.5) .. controls (8.3,1.6) and (7.6,2.7) .. (6.2,3.8);
  \draw[dashed,thick] (8.2,.5) .. controls (7.3,1.6) and (6.6,2.7) .. (5.2,3.8);
  \draw[dotted] (0,3.8)--(6.8,3.8);
  \node[left] at (0,3.8) {$t_c$};
  \fill (5.2,3.8) circle (1.8pt);
  \node[above right] at (5.2,3.8) {$(r_c,t_c)$};
  \node[align=left,anchor=west] at (.45,1.9)
    {$R_c'(t)=-\dfrac{q-1}{R_c(t)}$};
  \draw[<->] (7.2,-.42)--(9.2,-.42);
  \node[below] at (8.2,-.42) {$4h_c$};
  \node[anchor=west] at (7.75,2.7) {$h_c=r_c/\sqrt q$};
\end{tikzpicture}

\caption{Schematic image of the fixed cylinder \(P_2\) under
\eqref{eq:moving-cylinder}. Its radial width is \(4h_c\) and its time
length is \(4h_c^2\). Moving with \(R_c'(t)=-(q-1)/R_c(t)\) leaves
a bounded drift after rescaling. The drawing is not to scale.}
\label{fig:drift}
\end{figure}

Since \(R_c'=-(q-1)/R_c\), the function \(H=q(v-1)\), pulled back to
this cylinder, satisfies Lemma~\ref{lem:heat-perturbation} with
\begin{equation}\label{eq:H-coefficients}
\begin{aligned}
a&=z,\\
b&=h_c\left[(q-1)\left(\frac1r-\frac1{R_c}\right)-\frac zr\right],\\
F&=qh_c^2\left[\frac{(q-1)v(1-z)}{r^2}-pv^3u_r^2\right].
\end{aligned}
\end{equation}
Indeed \(|H|\le C_0\), since
\(q|v-1|=q|z-1|/(v+1)\), and \(|a-1|\le C_0/q\).
The apparently large drift is bounded by
\[
\left|h_c(q-1)\left(\frac1r-\frac1{R_c}\right)\right|
=\frac{(q-1)h_c^2|y|}{rR_c}\le C.
\]
For the source use \(qh_c^2=r_c^2\),
\((q-1)|1-z|\le C_0\), and \(r|u_r|\le2\).
All these coefficients and amplitudes are therefore bounded by
one numerical constant. Increasing \(q\) meets its smallness hypothesis.
It follows that \(|H_y|\le C\) on \(P_1\), and hence
\begin{equation}\label{eq:zr-cylinder}
|z_r|=\frac{2v}{qh_c}|H_y|\le\frac{C}{r_c\sqrt q}
\quad\text{on }P_1.
\end{equation}

Differentiate the first equation of~\eqref{eq:sideways}. With \(w=u_r\),
\begin{equation}\label{eq:ur-equation}
w_t=zw_{rr}+\left[z_r+\frac{q-1+z}{r}\right]w_r
+\left[\frac{z_r}{r}-\frac{q-1+z}{r^2}
          +2(p-1)e^{-2u}\right]w.
\end{equation}
On the smaller cylinder \(P_1\), set \(W=r_cw\). Its transformed drift is
\[
b_W=h_c(q-1)(1/r-1/R_c)+h_c(z_r+z/r),
\]
and its zeroth-order coefficient is
\[
c_W=h_c^2\left[z_r/r-(q-1+z)/r^2+2(p-1)e^{-2u}\right].
\]
Equations~\eqref{eq:uniform}, \eqref{eq:cylinder-inclusion}, and
\eqref{eq:zr-cylinder} bound \(|W|\), \(|b_W|\), and \(|c_W|\)
numerically. In particular \(h_c^2(p-1)e^{-2u}\le C\).
Treat \(c_WW\) as a bounded source and apply a fixed rescaled version of
Lemma~\ref{lem:heat-perturbation} on \(P_1\). Evaluating at its center
gives
\[
|W_y(0,0)|=r_ch_c|u_{rr}(r_c,t_c)|\le C.
\]
This proves the second estimate. The first follows from
\eqref{eq:zr-cylinder} at the center. Only the numerical \(C_0\) and
the one-dimensional estimates entered the constants, so neither
\(q\) nor \(k\) occurs in \(C\).
\end{proof}

\subsection{The Ricci norm and its reaction matrix}

The derivative estimates now control the radial sectional curvatures.
The remaining task is algebraic: write the curvature action on invariant
diagonal tensors in coordinates that include their multiplicities in the
norm. This isolates the leading term in \(q\).

The tensor evolution and the evolving metric give
\[
(\dt-\Delta)|\Ric|^2=-2|\nabla\Ric|^2
                         +4\langle\Rm(\Ric),\Ric\rangle.
\]
The inverse-metric derivatives in the squared norm cancel the cubic
Ricci term from the covariant tensor evolution. Kato's inequality then
gives, wherever \(|\Ric|>0\),
\begin{equation}\label{eq:ricci-norm}
(\dt-\Delta)|\Ric|
\le \frac{2}{|\Ric|}\langle\Rm(\Ric),\Ric\rangle.
\end{equation}
For diagonal tensors, our curvature convention is
\(\langle\Rm(D),D\rangle=\sum_{i,j}R_{ijij}D_iD_j\), with
\(R_{ijij}\) the sectional curvature for \(i\ne j\).

\begin{proposition}\label{prop:reaction}
At late times on \(0<r\le\frac18\sqrt\delta\), the reaction quotient
\[
c(x,t)=\frac{2\langle\Rm(\Ric),\Ric\rangle}{|\Ric|^2}
\quad\text{where }|\Ric|>0
\]
satisfies both bounds
\begin{equation}\label{eq:reaction-two}
c\le\frac{2(q-1)+C\sqrt q}{r^2},\qquad
c\le\frac{C_*}{\theta^2}.
\end{equation}
Here \(C\) is independent of \(q,k\); \(C_*\) may depend on the entire
chosen flow. The second bound holds at the pole orbits as well.
Globally at positive \(r\), one also has \(c\le D_1/r^2\) for some
flow-dependent \(D_1>0\).
\end{proposition}

\begin{proof}
The Ricci tensor is diagonal and scalar on each sphere factor. For a tensor
with entries \((\alpha,\beta^{\times p},\chi^{\times q})\), use the
Euclidean coordinates \((\alpha,\sqrt p\beta,\sqrt q\chi)\), so its
norm is the ordinary Euclidean norm. The curvature form has matrix
\begin{equation}\label{eq:reaction-matrix}
\begin{pmatrix}
0&\sqrt p\,h&\sqrt q\,h_q\\
\sqrt p\,h&(p-1)j&\sqrt{pq}\,\mu\\
\sqrt q\,h_q&\sqrt{pq}\,\mu&(q-1)\ell
\end{pmatrix}.
\end{equation}
Indeed its mixed terms are \(2ph\alpha\beta\),
\(2qh_q\alpha\chi\), and \(2pq\mu\beta\chi\); this accounts for
the ordered pairs and all multiplicities.

Using \(r\) as the spatial coordinate,
\begin{equation}\label{eq:curv-sideways}
\begin{gathered}
h=-z(u_{rr}+u_r^2)-\tfrac12z_ru_r,
\qquad h_q=-z_r/(2r),\qquad \mu=-zu_r/r,\\
j=e^{-2u}-zu_r^2,\qquad \ell=(1-z)/r^2.
\end{gathered}
\end{equation}
By Lemmas~\ref{lem:uniform} and~\ref{lem:largeq-derivatives}, after
multiplication by \(r^2\) the diagonal entries of
\eqref{eq:reaction-matrix} are bounded above by
\(0,q-1+C_0,C_0\), respectively. Its off-diagonal entries have sizes
\(O(\sqrt q),O(1),O(\sqrt q)\), respectively, with numerical constants.
The largest eigenvalue is therefore at most \(q-1+C\sqrt q\).
Equation~\eqref{eq:ricci-norm} proves the first bound.
The curvature contraction is bounded in absolute value by a
dimension-dependent constant times \(|\Rm||\Ric|^2\).
Proposition~\ref{prop:full-curvature} gives the other two bounds, including
the bound at a pole by continuity of the geometric curvature tensor.
\end{proof}
